% Plantilla: puente_losa.py reemplaza cada marcador (clave entre << y >>) por su valor calculado.
% No editar taller01_puentes.tex a mano: se regenera al ejecutar el script.
\documentclass[11pt,a4paper]{article}
\usepackage[utf8]{inputenc}
\usepackage[T1]{fontenc}
\usepackage{lmodern}
\usepackage[spanish,es-nodecimaldot,es-noshorthands,es-tabla]{babel}
\usepackage[a4paper,margin=2.1cm,headheight=14pt]{geometry}
\usepackage{amsmath,amssymb}
\usepackage{booktabs,array,tabularx}
\usepackage[table]{xcolor}
\usepackage{tikz}
\usetikzlibrary{arrows.meta,calc,patterns,decorations.pathreplacing}
\usepackage[most]{tcolorbox}
\usepackage{enumitem}
\usepackage{fancyhdr}
\usepackage{caption}
\usepackage{float}
\usepackage[hidelinks]{hyperref}

\definecolor{acento}{HTML}{1F4E79}
\definecolor{ok}{HTML}{1B7F3B}
\definecolor{mal}{HTML}{B3261E}
\definecolor{carga}{HTML}{2457C5}
\definecolor{nota}{HTML}{8A5A00}

\captionsetup{font=small,labelfont={bf,color=acento}}
\setlength{\parindent}{0pt}
\setlength{\parskip}{5pt}
\renewcommand{\arraystretch}{1.18}
\allowdisplaybreaks

\pagestyle{fancy}\fancyhf{}
\lhead{\small\color{acento} Taller de Puentes}
\rhead{\small Puente losa $L = 20$ ft \textperiodcentered{} Teorema de Barré}
\cfoot{\small\thepage}

\newtcolorbox{resultado}{colback=acento!6,colframe=acento,boxrule=0.8pt,arc=2pt,
  left=6pt,right=6pt,top=4pt,bottom=4pt}
\newtcolorbox{supuesto}{colback=nota!7,colframe=nota!70,boxrule=0.6pt,arc=2pt,
  left=6pt,right=6pt,top=3pt,bottom=3pt,fontupper=\small}

\newcommand{\uT}{\,\mathrm{t}}
\newcommand{\uTm}{\,\mathrm{t{\cdot}m}}
\newcommand{\uTmm}{\,\mathrm{t{\cdot}m/m}}
\newcommand{\uTpm}{\,\mathrm{t/m}}
\newcommand{\uKg}{\,\mathrm{kg/cm^2}}
\newcommand{\uCm}{\,\mathrm{cm}}
\newcommand{\uCmm}{\,\mathrm{cm^2/m}}
\newcommand{\uM}{\,\mathrm{m}}
\newcommand{\uMm}{\,\mathrm{mm}}
\newcommand{\ok}{\ensuremath{\;\textcolor{ok}{\checkmark}}}
\newcommand{\refn}[1]{\tag*{\small\color{acento}[#1]}}

% Viga simplemente apoyada de luz #1 (en las unidades del tikzpicture)
\newcommand{\vigaSA}[1]{%
  \draw[line width=1.8pt] (0,0) -- (#1,0);
  \draw[fill=white] (0,0) -- (-0.20,-0.36) -- (0.20,-0.36) -- cycle;
  \fill[pattern=north east lines] (-0.30,-0.36) rectangle (0.30,-0.48);
  \draw[fill=white] (#1,0) -- ({#1-0.20},-0.36) -- ({#1+0.20},-0.36) -- cycle;
  \draw ({#1-0.30},-0.44) -- ({#1+0.30},-0.44);
  \node[font=\scriptsize] at (0,-0.68) {A};
  \node[font=\scriptsize] at (#1,-0.68) {B};
}

\begin{document}

\begin{center}
{\LARGE\bfseries\color{acento} TALLER --- DISEÑO DE PUENTES}\\[6pt]
{\large \textbf{Pregunta 01:} puente losa de concreto armado, $L = 20$ ft}\\[2pt]
{\large \textbf{Pregunta 02:} reacción y momento máximo por el Teorema de Barré}\\[10pt]
\begin{tabular}{>{\bfseries}r p{12.9cm}}
Alumno: & \rule{8cm}{0.4pt}\\
Configuración: & A--L: $f'_c = 420\uKg$, asfalto 3'', ejes HL-93 $d_{1-2} = 3.00\uM$ y $d_{2-3} = 3.50\uM$\\
Normas: & AASHTO LRFD Bridge Design Specifications (base del Manual de Puentes MTC-2018);
          Manual de Carreteras: Diseño Geométrico DG-2018\\
Unidades: & cargas en t y m; diseño en kg y cm (las fórmulas AASHTO en mm o MPa se evalúan en esas unidades)
\end{tabular}
\end{center}

%==========================================================================
\section*{0. Criterios de clase aplicados}

\begin{enumerate}[leftmargin=*,label=\textbf{\arabic*.},itemsep=3pt]
\item \textbf{Estados límite.} Servicio (S), Fatiga y Fractura (FxF), Resistencia (R) y Evento Extremo (EE).
  El momento de diseño $M_u$ sale de \textbf{Resistencia I}; con Servicio I se controla la fisuración y con
  Fatiga I el rango de esfuerzos del acero. Evento Extremo (sismo, socavación, choque) no gobierna la flexión
  de una losa simplemente apoyada: se revisa en apoyos y subestructura.
\item \textbf{Tipo de estructura según la luz.} Para $L = 50$ m no sirven ni la viga ni la losa de concreto
  armado: se pasa a pretensado o postensado, o se acorta la luz con apoyos intermedios (25\,+\,25 m;
  20\,+\,10\,+\,20 m). Para $L = 6.096$ m la solución natural es una \textbf{losa maciza simplemente apoyada}.
\item \textbf{Peraltes mínimos} (AASHTO Tabla 2.5.2.6.3-1, $S$ = luz en mm). Los valores de la pizarra
  salen de estas fórmulas:

  \begin{center}\small
  \begin{tabular}{lcc}
  \toprule
  Superestructura & Simplemente apoyada (SA) & Continua (C)\\
  \midrule
  Losa, refuerzo paralelo al tráfico & $t = \dfrac{1.2\,(S+3000)}{30}$ & $t = \dfrac{S+3000}{30} \ge 165\uMm$\\[10pt]
  Viga T de concreto armado & $h = 0.070\,L$ & $h = 0.065\,L$\\
  \bottomrule
  \end{tabular}
  \end{center}
  Pizarra: SA, $L = 20$ m $\rightarrow$ $t = 1.2(23\,000)/30 = 920\uMm \approx 0.90\uM$ y $h = 0.070(20) = 1.40\uM$;
  C, $L = 50$ m $\rightarrow$ $t = 53\,000/30 = 1767\uMm \approx 1.77\uM$ y $h = 0.065(50) = 3.25\uM$.
\item \textbf{Carga viva HL-93} (AASHTO 3.6.1.2): camión de $3.63$--$14.52$--$14.52\uT$ (separaciones 4.27 m y
  de 4.27 a 9.00 m), tándem de $2\times 11.34\uT$ separados 1.20 m y carga de carril de $0.952\uTpm$ en 3.00 m
  de ancho. Se toma el mayor efecto entre camión y tándem y se le suma el carril. En este taller las
  separaciones del camión cambian según el apellido: $d_{1-2} = 3.00\uM$ y $d_{2-3} = 3.50\uM$.
\end{enumerate}

%==========================================================================
\section*{PREGUNTA 01 --- Puente losa de concreto armado, $L = 20$ ft}

\subsection*{1.1 Datos}
\begin{center}\small
\begin{tabular}{ll@{\hspace{1.5cm}}ll}
\toprule
Luz & $L = 20\,\text{ft}\times 0.3048 = 6.096\uM$ & Concreto & $f'_c = 420\uKg$ ($41.19$ MPa)\\
Tránsito & IMDA $= 4\,500$ veh/día & Acero & $f_y = 4\,200\uKg$ (Grado 60)\\
Orografía & terreno accidentado (tipo 3) & Asfalto & $3'' = 7.62\uCm$\\
Velocidad de diseño & $V = 70$ km/h & Ejes HL-93 & $d_{1-2} = 3.00\uM$, $d_{2-3} = 3.50\uM$\\
Gálibo hidráulico & $\ge 2.50\uM$ & Importancia & puente importante\\
\bottomrule
\end{tabular}
\end{center}

\subsection*{1.2 Clasificación de la vía y sección transversal}
Con IMDA $= 4\,500$ veh/día $\in [4\,001;\ 6\,000]$ la vía es \textbf{Autopista de Segunda Clase} (DG-2018, Art.~101.01):
calzadas separadas por un separador central de 1.00 a 6.00 m, cada una con dos o más carriles de 3.60 m.
Como las calzadas van separadas, se proyectan \textbf{dos puentes gemelos}, uno por sentido, iguales entre sí.
Anchos de cada puente, según las tablas de clase (DG-2018) para autopista de 2.ª clase, orografía 3 y
$V = 70$ km/h:
\begin{itemize}[leftmargin=*,itemsep=1pt]
\item \textbf{Calzada:} 7.20 m, 2 carriles de 3.60 m (Tabla 304.01).
\item \textbf{Bermas:} 3.00 m a cada lado (Tabla 304.02).
\item \textbf{Vereda} (criterio de clase, ``frontera de la velocidad''): con $V < 70$ km/h basta una vereda de 1.20 m;
  con $V \ge 70$ km/h la vereda efectiva mínima es 1.50 m y se separa del tránsito con una barrera rígida tipo
  New Jersey, con baranda peatonal exterior de altura $\ge 1.10$ m. Como $V = 70 \ge 70$ km/h
  $\Rightarrow$ \textbf{vereda de 1.50 m}, de concreto $f'_c = 280$ kg/cm$^2$ y 0.20 m de espesor
  (sardinel de 150 a 200 mm).
\item \textbf{Baranda metálica peatonal:} 0.15 m de ancho y 0.15 t/m (valor adoptado en proyectos
  peruanos: 0.10 a 0.15 t/m).
\item \textbf{Barrera New Jersey}: $H = 0.813$ m, ancho $B = 0.40$ m; para el peso se usa el área de la
  tabla de clase, $0.2540\ \mathrm{m^2}$ (conservador).
\end{itemize}
La vereda va en el lado exterior de cada puente; hacia el separador central basta la barrera.

\begin{figure}[H]
\centering
\begin{tikzpicture}[x=1.0cm,y=1.0cm]
\fill[gray!25] (2.050,0) rectangle (15.250,0.0762);
\draw[fill=gray!8] (0,-0.400) rectangle (15.650,0);
\node[font=\scriptsize] at (7.825,-0.200) {Losa maciza $t = 0.40$ m};
\draw[fill=gray!15] (0.150,0) rectangle (1.650,0.200);
\draw[fill=gray!55] (0,0) rectangle (0.150,1.300);
\node[font=\tiny, rotate=90, anchor=south] at (0,0.750) {Baranda};
\draw[fill=gray!35] (1.650,0.000) -- (2.050,0.000) -- (2.050,0.076) -- (1.817,0.330) -- (1.750,0.813) -- (1.650,0.813) -- cycle;
\draw[fill=gray!35] (15.650,0.000) -- (15.250,0.000) -- (15.250,0.076) -- (15.483,0.330) -- (15.550,0.813) -- (15.650,0.813) -- cycle;
\draw[<->, >=Stealth] (0.000,1.55) -- (0.150,1.55) node[midway, above, font=\tiny] {0.15};
\draw[gray] (0.000,1.45) -- (0.000,1.6500000000000001);
\draw[<->, >=Stealth] (0.150,1.55) -- (1.650,1.55) node[midway, above, font=\scriptsize] {1.50};
\draw[gray] (0.150,1.45) -- (0.150,1.6500000000000001);
\node[font=\scriptsize] at (0.900,0.350) {Vereda};
\draw[<->, >=Stealth] (1.650,1.55) -- (2.050,1.55) node[midway, above, font=\tiny] {0.40};
\draw[gray] (1.650,1.45) -- (1.650,1.6500000000000001);
\draw[<->, >=Stealth] (2.050,1.55) -- (5.050,1.55) node[midway, above, font=\scriptsize] {3.00};
\draw[gray] (2.050,1.45) -- (2.050,1.6500000000000001);
\node[font=\scriptsize] at (3.550,0.350) {Berma ext.};
\draw[<->, >=Stealth] (5.050,1.55) -- (8.650,1.55) node[midway, above, font=\scriptsize] {3.60};
\draw[gray] (5.050,1.45) -- (5.050,1.6500000000000001);
\node[font=\scriptsize] at (6.850,0.350) {Carril 1};
\draw[<->, >=Stealth] (8.650,1.55) -- (12.250,1.55) node[midway, above, font=\scriptsize] {3.60};
\draw[gray] (8.650,1.45) -- (8.650,1.6500000000000001);
\node[font=\scriptsize] at (10.450,0.350) {Carril 2};
\draw[<->, >=Stealth] (12.250,1.55) -- (15.250,1.55) node[midway, above, font=\scriptsize] {3.00};
\draw[gray] (12.250,1.45) -- (12.250,1.6500000000000001);
\node[font=\scriptsize] at (13.750,0.350) {Berma int.};
\draw[<->, >=Stealth] (15.250,1.55) -- (15.650,1.55) node[midway, above, font=\tiny] {0.40};
\draw[gray] (15.250,1.45) -- (15.250,1.6500000000000001);
\draw[gray] (15.650,1.45) -- (15.650,1.6500000000000001);
\draw[dashed] (8.650,0.0762) -- (8.650,0.20);
\draw[<->, >=Stealth] (0,-0.850) -- (15.650,-0.850) node[midway, fill=white, font=\scriptsize, inner sep=1pt] {$W = 15.65$ m};
\draw[decorate, decoration={brace, mirror, amplitude=4pt}] (0,-0.480) -- (1.636,-0.480) node[midway, below=3pt, font=\scriptsize] {$E_{b,v}$};
\draw[decorate, decoration={brace, amplitude=4pt}] (15.650,-0.480) -- (14.132,-0.480) node[midway, below=3pt, font=\scriptsize] {$E_{b,s}$};
\end{tikzpicture}
\caption{Sección transversal de un puente (un sentido). Cotas en m. $E_{b,v}$ y $E_{b,s}$: franjas de borde
del lado de la vereda y del lado del separador.}
\end{figure}

Ancho de calzada, ancho total y número de carriles de diseño:
\begin{align*}
w &= 3.00 + 7.20 + 3.00 = 13.20\uM \quad\text{(entre caras de barreras)}\\
W &= \underbrace{0.15 + 1.50 + 0.40}_{\text{baranda + vereda + barrera}} + 13.20 + 0.40
 = \boxed{15.65\uM}\\
N_L &= \operatorname{ENT}\!\left(\frac{w}{3.60}\right) = \operatorname{ENT}\!\left(\frac{13.20}{3.60}\right)
 = \operatorname{ENT}(3.67) = \boxed{3\ \text{carriles de diseño}} \refn{AASHTO 3.6.1.1.1}
\end{align*}
Aunque solo se pintan 2 carriles, AASHTO cuenta el ancho de las bermas porque en el futuro pueden usarse como carril.


\subsection*{1.3 Tipo de superestructura y predimensionamiento}
Losa maciza de concreto armado, simplemente apoyada, con el refuerzo principal paralelo al tráfico.
\begin{align*}
t_{\min} &= \frac{1.2\,(S+3000)}{30} = \frac{1.2\,(6\,096+3000)}{30} = 363.8\uMm
 \refn{AASHTO Tabla 2.5.2.6.3-1}\\
&\Rightarrow\ \boxed{t = 0.40\uM} \ \ge t_{\min} = 0.364\uM \ok
\end{align*}

\subsection*{1.4 Materiales y parámetros de la sección}
Propiedades del concreto y del acero, y peralte efectivo con una capa de barras:
\begin{align*}
E_c &= 15\,344\sqrt{f'_c} = 15\,344\sqrt{420} = 314\,460\uKg
 \refn{AASHTO 5.4.2.4, $w_c = 2320$ kg/m$^3$}\\
E_s &= 200\,000\ \text{MPa} = 2\,039\,432\uKg \qquad
n = \frac{E_s}{E_c} = 6.49 \refn{AASHTO 5.4.3.2}\\
\beta_1 &= 0.85 - 0.05\,\frac{f'_c-28}{7} = 0.85 - 0.05\,\frac{41.19-28}{7} = 0.756
 \quad(f'_c \text{ en MPa})
 \refn{AASHTO 5.6.2.2}\\
f_r &= 0.63\sqrt{f'_c\,[\text{MPa}]} = 0.63\sqrt{41.19}\ \text{MPa} = 41.23\uKg \refn{AASHTO 5.4.2.6}\\
d &= t - r_e = 40 - 4.5 = 35.50\uCm
 \quad(r_e = 4.5\uCm \text{ del fondo al eje de la barra}) \refn{criterio de clase}
\end{align*}

\subsection*{1.5 Modificador de cargas}
$\eta_D = 1.00$ (diseño convencional), $\eta_R = 1.00$ (redundancia convencional) y $\eta_I = 1.05$
(\textbf{puente importante}, dato del problema):
\begin{align*}
\eta = \eta_D\,\eta_R\,\eta_I = 1.00 \times 1.00 \times 1.05 = 1.05 \refn{AASHTO 1.3.2 a 1.3.5}
\end{align*}
En Servicio y Fatiga se usa $\eta = 1.00$.

\subsection*{1.6 Cargas permanentes (franja interior de 1 m)}
Pesos específicos de la tabla de clase (MTC-2016): con $f'_c = 420 > 350$ kg/cm$^2$ el concreto simple pesa
$2240 + 0.23\,f'_c$ y el armado 80 kg/m$^3$ más; la superficie bituminosa pesa 2240 kg/m$^3$.
\begin{align*}
\gamma_c &= 2240 + 0.23(420) + 80 = 2\,416.6\ \mathrm{kg/m^3} \qquad
\gamma_{c,\text{vereda}} = 2320 + 80 = 2400\ \mathrm{kg/m^3}\ \ (f'_c = 280) \refn{MTC-2016}\\
w_{DC} &= t\,\gamma_c = 0.40 \times 2.4166 = 0.967\uTpm\\
w_{DW} &= e\,\gamma_{asf} = 0.0762 \times 2.24 = 0.1707\uTpm\\
M_{DC} &= w_{DC}\,\frac{L^2}{8} = 0.967 \times \frac{6.096^2}{8} = 0.967 \times 4.6452 = 4.490\uTmm\\
M_{DW} &= w_{DW}\,\frac{L^2}{8} = 0.1707 \times 4.6452 = 0.793\uTmm
\end{align*}

\subsection*{1.7 Carga viva vehicular por carril}

\paragraph{Teorema de Barré.} En una viga simplemente apoyada recorrida por un tren de cargas, el momento máximo
absoluto ocurre bajo una de las cargas, cuando el centro de la luz divide en partes iguales la distancia $e$ entre
esa carga y la resultante $R$ de las cargas que están sobre la viga. La carga crítica queda a $x = L/2 + e/2$
del apoyo y
\begin{gather*}
M_{\max} = \frac{R}{L}\left(\frac{L}{2} - \frac{|e|}{2}\right)^{2} - \sum_{\text{lado opuesto a } R} P_j\,|x_j - x_k|\\
\text{Dos cargas iguales } P \text{ separadas } s \ (e = s/2):\qquad
M_{\max} = \frac{2P}{L}\left(\frac{L}{2} - \frac{s}{4}\right)^{2}
\end{gather*}
La posición solo vale si en ella ningún otro eje entra o sale de la luz. Si no, el máximo de ese tramo
ocurre cuando un eje está sobre un apoyo, y se evalúa ahí.

\paragraph{a) Camión HL-93 modificado} ($3.63$ --- $3.00$ m --- $14.52$ --- $3.50$ m --- $14.52$ t).
Su largo es $3.00 + 3.50 = 6.50\uM > L = 6.096\uM$: \textbf{los tres ejes no caben a la vez},
así que se analizan los grupos de dos ejes y el eje solo.

\begin{center}\small
\begin{tabular}{ccccccc}
\toprule
Grupo & Eje crítico & $R$ (t) & $e$ (m) & $x$ crítico (m) & ¿Posición válida? & $M$ (t$\cdot$m)\\
\midrule
1--2 & 1 & 18.15 & -2.400 & 1.848 & \textcolor{ok}{sí} & 10.168 \\
1--2 & 2 & 18.15 & 0.600 & 3.348 & \textcolor{ok}{sí} & 22.484 \\
2--3 & 2 & 29.04 & -1.750 & 2.173 & \textcolor{ok}{sí} & 22.494 \\
2--3 & 3 & 29.04 & 1.750 & 3.923 & \textcolor{ok}{sí} & 22.494 \\
\bottomrule
\end{tabular}
\end{center}
Un eje de 14.52 t solo al centro: $M = PL/4 = 14.52\times6.096/4 = 22.128\uTm$. Gobierna el grupo de ejes
2, 3:
\begin{equation*}
M_{\text{camión}} = \frac{2(14.52)}{6.096}\left(\frac{6.096}{2} - \frac{3.50}{4}\right)^{2}
 = 22.494\uTm \quad\text{(bajo el eje 2, a } x = 2.173\uM\text{)}
\end{equation*}

\begin{figure}[H]
\centering
\begin{tikzpicture}[x=1.7224cm,y=1cm]
\clip (-1.4,-1.65) rectangle ({6.0960+1.4},2.15);
\vigaSA{6.0960}
\draw[-{Stealth[length=2.2mm]}, very thick, gray!55] (-0.827,1.25) -- (-0.827,0.10) node[pos=0, above, font=\scriptsize, text=gray!55, align=center] {\textbf{1}\\3.63};
\draw[-{Stealth[length=2.2mm]}, very thick, carga] (2.173,1.25) -- (2.173,0.10) node[pos=0, above, font=\scriptsize, text=carga, align=center] {\textbf{2}\\14.52};
\draw[-{Stealth[length=2.2mm]}, very thick, carga] (5.673,1.25) -- (5.673,0.10) node[pos=0, above, font=\scriptsize, text=carga, align=center] {\textbf{3}\\14.52};
\draw[dashed,gray] ({6.0960/2},-0.25) -- ({6.0960/2},1.95) node[above,font=\scriptsize]{$L/2$};
\draw[<->,>=Stealth] (0,-1.0) -- (2.1730,-1.0) node[midway,below,font=\scriptsize]{$x = 2.173$ m};
\draw[<->,>=Stealth] (0,-1.45) -- (6.0960,-1.45) node[midway,fill=white,font=\scriptsize,inner sep=1pt]{$L = 6.096$ m};
\end{tikzpicture}
\caption{Camión HL-93 modificado en la posición de Barré (en gris, el eje que queda fuera de la luz).}
\end{figure}

\paragraph{b) Tándem} ($2 \times 11.34\uT$ a 1.20 m):
\begin{equation*}
M_{\text{tándem}} = \frac{2(11.34)}{6.096}\left(\frac{6.096}{2} - \frac{1.20}{4}\right)^{2}
 = \frac{22.68}{6.096}\,(2.748)^{2} = 28.095\uTm \quad(x = 3.348\uM)
\end{equation*}

\begin{figure}[H]
\centering
\begin{tikzpicture}[x=1.7224cm,y=1cm]
\clip (-1.4,-1.65) rectangle ({6.0960+1.4},2.15);
\vigaSA{6.0960}
\draw[-{Stealth[length=2.2mm]}, very thick, carga] (2.148,1.25) -- (2.148,0.10) node[pos=0, above, font=\scriptsize, text=carga, align=center] {\textbf{1}\\11.34};
\draw[-{Stealth[length=2.2mm]}, very thick, carga] (3.348,1.25) -- (3.348,0.10) node[pos=0, above, font=\scriptsize, text=carga, align=center] {\textbf{2}\\11.34};
\draw[dashed,gray] ({6.0960/2},-0.25) -- ({6.0960/2},1.95) node[above,font=\scriptsize]{$L/2$};
\draw[<->,>=Stealth] (0,-1.0) -- (3.3480,-1.0) node[midway,below,font=\scriptsize]{$x = 3.348$ m};
\draw[<->,>=Stealth] (0,-1.45) -- (6.0960,-1.45) node[midway,fill=white,font=\scriptsize,inner sep=1pt]{$L = 6.096$ m};
\end{tikzpicture}
\caption{Tándem de diseño en la posición de Barré.}
\end{figure}

\paragraph{c) Carga de carril:} $M_{\text{carril}} = 0.952 \times 4.6452 = 4.422\uTm$.

\paragraph{d) Momento por carril con impacto.} $IM = 33\%$ solo para camión o tándem, no para el carril
(AASHTO 3.6.2.1). Gobierna el \textbf{tándem} ($28.095 > 22.494$):
\begin{align*}
M_{LL+IM} &= (1+IM)\,M_{\text{tándem}} + M_{\text{carril}} = 1.33 \times 28.095 + 4.422\\
 &= 37.367 + 4.422 = \boxed{41.789\uTm \text{ por carril}}
\end{align*}

\subsection*{1.8 Ancho de franja equivalente}
Fórmulas de clase (AASHTO 4.6.2.3, en metros), con $L_1 = \min(L;\ 18.26) = 6.096\uM$ y el ancho total
$W = 15.65\uM$:
\begin{align*}
\intertext{\textbf{Un solo carril cargado}, $W_1 = \min(W;\ 9.14) = 9.14\uM$:}
E_1 &= 0.25 + 0.42\sqrt{L_1 W_1} = 0.25 + 0.42\sqrt{6.096\times9.14} = 0.25 + 0.42\,(7.4644) = \boxed{3.385\uM}
\intertext{\textbf{Múltiples carriles cargados}, $W_1 = \min(W;\ 18) = 15.65\uM$:}
E_m &= 2.10 + 0.12\sqrt{L_1 W_1} = 2.10 + 0.12\sqrt{6.096\times15.65} = 2.10 + 0.12\,(9.7674)
 = 3.272\uM\\
 &\le \frac{W}{N_L} = \frac{15.65}{3} = 5.217\uM \quad\Rightarrow\quad E_m = \boxed{3.272\uM}
\end{align*}

\textbf{Momento de sobrecarga por metro} (fórmula de clase $M_{LL} = m\,M_{SC\text{-}V}/E \times b$, con $b = 1$ m y
$m$ de la Tabla 3.6.1.1.2-1): se calcula con un carril y con varios carriles, y se toma el mayor.
\begin{align*}
\text{Un carril } (m = 1.20):\quad M_{LL+IM} &= \frac{1.20 \times 41.789}{3.385} = 14.814\uTmm\\
\text{Múltiples carriles } (m = 1.00):\quad M_{LL+IM} &= \frac{1.00 \times 41.789}{3.272} = 12.771\uTmm\\
\Rightarrow\quad M_{LL+IM} &= \max(14.814;\ 12.771) = \boxed{14.814\uTmm}
\end{align*}
\textbf{Gobierna un carril.}
{\small Para múltiples carriles se usa $m = 1.00$ (2 carriles cargados, como en clase); con 3 carriles,
$m = 0.85$ daría un momento menor.}


\textbf{Franjas de borde} (AASHTO 4.6.2.1.4b): distancia del borde a la cara de la barrera + 0.30 m + $E/4$, sin
pasar de $E/2$ ni de 1.80 m, con $E = \min(E_1;\ E_m) = 3.272$ m:
\begin{align*}
\text{Lado separador:}\quad E_{b,s} &= 0.40 + 0.30 + \frac{3.272}{4} = 1.518\uM
 \ \le\ \min(1.636;\ 1.80) \ \Rightarrow\ \boxed{E_{b,s} = 1.518\uM}\\
\text{Lado vereda:}\quad E_{b,v} &= 2.05 + 0.30 + \frac{3.272}{4} = 3.168\uM
 \ \le\ \min(1.636;\ 1.80) \ \Rightarrow\ \boxed{E_{b,v} = 1.636\uM}
\end{align*}
Del lado de la vereda la franja ($1.636$ m) no llega a la cara de la barrera, que está a $2.05$ m del borde:
ningún vehículo la carga. Lleva losa, vereda, barrera, baranda y la carga peatonal.

\subsection*{1.9 Momentos por metro de ancho y combinaciones de carga}
\textbf{Franja interior:} $M_{LL+IM} = 14.814\uTmm$ (mayor entre un carril y múltiples carriles, sección 1.8).

\textbf{Franja de borde, lado separador} ($E_{b,s} = 1.518$ m): la barrera
($0.2540\ \mathrm{m^2}\times 2.40 = 0.610\uTpm$) va en la franja; el asfalto solo cubre
$E_{b,s} - 0.40 = 1.118\uM$; la carga viva es una línea de ruedas más la parte del carril que
cae en la franja, con $m = 1.20$ (un carril cargado):
\begin{align*}
w_{DC} &= 0.967 + \frac{0.610}{1.518} = 1.368\uTpm \qquad
w_{DW} = 0.1707\times\frac{1.118}{1.518} = 0.1257\uTpm\\
M_{LL+IM} &= \frac{m\left[\,0.5\,(1.33)\,M_{\text{tándem}} + M_{\text{carril}}\,\dfrac{E_{b,s} - 0.40}{3.00}\right]}{E_{b,s}}
 = \frac{1.20\left[\,0.5(1.33)(28.095) + 4.422\,\dfrac{1.118}{3.00}\right]}{1.518}\\
 &= \frac{1.20\,(20.331)}{1.518} = 16.072\uTmm
\end{align*}

\textbf{Franja de borde, lado vereda} ($E_{b,v} = 1.636$ m): vereda ($0.20 \times 2.40 \times 1.50
= 0.720\uTpm$), barrera ($0.610\uTpm$), baranda ($0.15\uTpm$) y carga peatonal
$PL = 3.6\ \text{kPa} = 0.367\ \mathrm{t/m^2}$ en toda la vereda (AASHTO 3.6.1.6):
\begin{align*}
w_{DC} &= 0.967 + \frac{0.720 + 0.610 + 0.15}{1.636} = 1.871\uTpm \qquad
w_{PL} = \frac{0.367 \times 1.50}{1.636} = 0.336\uTpm\\
M_{DC} &= 1.871 \times 4.6452 = 8.691\uTmm \qquad M_{PL} = 0.336 \times 4.6452 = 1.563\uTmm
\end{align*}

\textbf{Combinaciones} (AASHTO Tabla 3.4.1-1):
\begin{align*}
\text{Resistencia I:}\quad M_u &= \eta\,\big[1.25\,M_{DC} + 1.50\,M_{DW} + 1.75\,(M_{LL+IM} + M_{PL})\big]\\
\text{Servicio I:}\quad M_s &= 1.00\,M_{DC} + 1.00\,M_{DW} + 1.00\,(M_{LL+IM} + M_{PL})
\end{align*}
Reemplazando los momentos de cada franja:
\begin{align*}
M_u^{\text{int}} &= 1.05\,\big[1.25(4.490) + 1.50(0.793) + 1.75(14.814)\big] = 1.05\,(32.727)
 = \boxed{34.363\uTmm}\\
M_u^{\text{sep.}} &= 1.05\,\big[1.25(6.356) + 1.50(0.584) + 1.75(16.072)\big] = 1.05\,(36.946)
 = \boxed{38.794\uTmm}\\
M_u^{\text{vereda}} &= 1.05\,\big[1.25(8.691) + 1.75(1.563)\big] = 1.05\,(13.599)
 = \boxed{14.279\uTmm}
\end{align*}

\begin{center}\small
\begin{tabular}{lcccccccc}
\toprule
Franja & Ancho (m) & $M_{DC}$ & $M_{DW}$ & $M_{LL+IM}$ & $M_{PL}$ & $M_u$ (Resist. I) & $M_s$ (Serv. I) \\
 & & \multicolumn{6}{c}{(t$\cdot$m/m)}\\
\midrule
Interior & 1.000 & 4.490 & 0.793 & 14.814 & -- & \textbf{34.363} & 20.097\\
Borde, lado separador & 1.518 & 6.356 & 0.584 & 16.072 & -- & \textbf{38.794} & 23.011\\
Borde, lado vereda & 1.636 & 8.691 & -- & -- & 1.563 & \textbf{14.279} & 10.254\\
\bottomrule
\end{tabular}
\end{center}

\subsection*{1.10 Diseño por flexión (Resistencia I)}
Sección rectangular $b = 100\uCm$, $d = 35.50\uCm$, $\phi = 0.90$:
\begin{equation*}
A_s = \frac{0.85\,f'_c\,b\,d}{f_y}\left[1 - \sqrt{1 - \frac{2\,M_u}{\phi\,0.85\,f'_c\,b\,d^{2}}}\,\right]
\qquad a = \frac{A_s f_y}{0.85\,f'_c\,b} \qquad \phi M_n = \phi A_s f_y\left(d - \frac{a}{2}\right)
\refn{AASHTO 5.6.3.2}
\end{equation*}

\textbf{Franja interior} ($M_u = 3\,436\,305$ kg$\cdot$cm):
\begin{align*}
\frac{2M_u}{\phi\,0.85 f'_c b d^2} &= 0.1697 \quad\Rightarrow\quad
A_s = \frac{0.85(420)(100)(35.50)}{4\,200}\left[1-\sqrt{1-0.1697}\right] = 26.80\uCmm\\
s &= \frac{100\,A_b}{A_s} = \frac{100(5.10)}{26.80} = 19.0\uCm
 \ \Rightarrow\ \boxed{\varnothing\,1'' \ @\ 17.5\uCm}\quad (A_s = 29.14\uCmm)\\
a &= \frac{29.14(4\,200)}{0.85(420)(100)} = 3.43\uCm \qquad
c = \frac{a}{\beta_1} = 4.54\uCm \qquad \frac{c}{d} = 0.128 \le 0.42 \ok\\
\phi M_n &= 0.90(29.14)(4\,200)\left(35.50 - \frac{3.43}{2}\right) = 37.218\uTmm
 \ \ge\ M_u = 34.363 \ok\ (D/C = 0.92)
\end{align*}

\textbf{Franja de borde, lado separador} ($M_u = 3\,879\,361$ kg$\cdot$cm):
\begin{align*}
\frac{2M_u}{\phi\,0.85 f'_c b d^2} &= 0.1916 \quad\Rightarrow\quad
A_s = \frac{0.85(420)(100)(35.50)}{4\,200}\left[1-\sqrt{1-0.1916}\right] = 30.45\uCmm\\
s &= \frac{100(5.10)}{30.45} = 16.8\uCm
 \ \Rightarrow\ \boxed{\varnothing\,1'' \ @\ 15.0\uCm}\quad (A_s = 34.00\uCmm)\\
a &= 4.00\uCm \qquad c = 5.29\uCm \qquad \frac{c}{d} = 0.149 \le 0.42 \ok\\
\phi M_n &= 0.90(34.00)(4\,200)\left(35.50 - \frac{4.00}{2}\right) = 43.054\uTmm
 \ \ge\ M_u = 38.794 \ok\ (D/C = 0.90)
\end{align*}

\textbf{Franja de borde, lado vereda} ($M_u = 1\,427\,917$ kg$\cdot$cm):
\begin{align*}
\frac{2M_u}{\phi\,0.85 f'_c b d^2} &= 0.0705 \quad\Rightarrow\quad A_s = 10.84\uCmm\\
s &= \frac{100(5.10)}{10.84} = 47.1\uCm \ \Rightarrow\ 45.0\uCm\ \text{(máximo 45 cm)}
\end{align*}
Por resistencia bastaría $\varnothing\,1''\ @\ 45.0$ cm, pero no cumple fisuración (Servicio I). Se cierra el
espaciamiento hasta cumplir: $\boxed{\varnothing\,1''\ @\ 30.0\uCm}$ ($A_s = 17.00\uCmm$),
$\phi M_n = 22.170\uTmm \ge 14.279$ \ok, $c/d = 0.075 \le 0.42$ \ok.

\subsection*{1.11 Verificaciones}

\paragraph{a) Acero mínimo} (AASHTO 5.6.3.3): $\phi M_n \ge \min(M_{cr};\ 1.33\,M_u)$, con
$\gamma_1 = 1.6$ y $\gamma_3 = 0.67$ (A615 Gr 60).
\begin{align*}
M_{cr} = \gamma_3\,\gamma_1\,f_r\,S_c = 0.67(1.6)(41.23)\,\frac{100(40)^2}{6}
 = 0.67(1.6)(41.23)(26\,667) = 11.786\uTmm
\end{align*}
Interior: $\min(11.786;\ 45.703) = 11.786 \le \phi M_n = 37.218$ \ok\qquad
Borde sep.: $11.786 \le 43.054$ \ok\qquad Borde vereda: $\min(11.786;\ 18.991) = 11.786 \le 22.170$ \ok

Cuantía mínima (criterio de clase):
\begin{equation*}
\rho_{\min} = 0.03\,\frac{f'_c}{f_y} = 0.03\,\frac{420}{4\,200} = 0.0030 \ \Rightarrow\
A_{s,\min} = 0.0030(100)(35.50) = 10.65\uCmm
\end{equation*}
$A_{s,\min} = 10.65 \le 29.14$ (interior), $34.00$ (borde separador) y $17.00$ cm$^2$/m (borde vereda) \ok

\paragraph{b) Fisuración, Servicio I} (AASHTO 5.6.7). Sección fisurada transformada, $n = 6.49$:
\begin{align*}
\frac{b\,y^2}{2} = n A_s (d - y) \qquad j d = d - \frac{y}{3} \qquad
f_{ss} = \frac{M_s}{A_s\,j d} \le 0.6 f_y = 2\,520\uKg\\
d_c = r_e = 4.50\uCm \qquad
\beta_s = 1 + \frac{d_c}{0.7\,(h - d_c)} = 1.1811 \qquad
s \le \frac{123\,000\,\gamma_e}{\beta_s\,f_{ss}} - 2 d_c \ \ [\text{mm, MPa}]
\end{align*}
Se usa $\gamma_e = 0.75$ (exposición clase 2, por la humedad del río; con clase 1, $\gamma_e = 1.00$,
cumple con más holgura).
\begin{center}\footnotesize
\begin{tabular}{lcccccccc}
\toprule
Franja & $A_s$ & $nA_s$ & $y$ & $jd$ & $M_s$ & $f_{ss}$ & $s_{\max}$ & $s$\\
 & (cm$^2$/m) & (cm$^2$) & (cm) & (cm) & (t$\cdot$m/m) & (kg/cm$^2$) & (cm) & (cm)\\
\midrule
Interior & 29.14 & 189.01 & 9.85 & 32.22 & 20.097 & 2\,140.5 & 28.2 & 17.5 \ok\\
Borde sep. & 34.00 & 220.51 & 10.50 & 32.00 & 23.011 & 2\,115.0 & 28.7 & 15.0 \ok\\
Borde vereda & 17.00 & 110.25 & 7.81 & 32.90 & 10.254 & 1\,833.6 & 34.4 & 30.0 \ok\\
\bottomrule
\end{tabular}
\end{center}
Ejemplo (interior): $f_{ss} = 2\,140.5\uKg = 209.9$ MPa,
$s_{\max} = \dfrac{123\,000(0.75)}{1.1811(209.9)} - 2(45.0)\ \text{mm} = 28.2\uCm$.

\paragraph{c) Fatiga I} (AASHTO 3.6.1.4 y 5.5.3). Un camión con $d_{1-2} = 3.00\uM$ y 9.00 m entre ejes
traseros, $IM = 15\%$, franja de un carril $E_1$ y sin presencia múltiple (se divide entre 1.20,
AASHTO 3.6.1.1.2). Con 9.00 m solo entran los ejes 1 y 2: $M = 22.484\uTm$ por carril.
\begin{align*}
M_{fat} &= \frac{(1+0.15)\,M}{1.20\,E_1} = \frac{1.15(22.484)}{1.20(3.385)} = 6.365\uTmm \qquad
\Delta f = \frac{\gamma\,M_{fat}}{A_s\,j d},\ \ \gamma = 1.50\\
f_{\min} &= \frac{M_{DC} + M_{DW}}{A_s\,j d} \qquad
\Delta F_{TH} = 24 - 0.33 f_{\min}\ [\text{ksi}] = 1687 - 0.33 f_{\min}\ [\text{kg/cm}^2]
\end{align*}
Interior: $f_{\min} = 562.7$, $\Delta f = 1\,016.9 \le \Delta F_{TH} = 1\,501.7\uKg$ \ok\qquad
Borde sep.: $f_{\min} = 637.8$, $\Delta f = 877.6 \le 1\,476.9\uKg$ \ok

La franja del lado de la vereda no recibe carga vehicular, así que no se revisa por fatiga.

{\small Con la 9.ª edición de AASHTO ($\gamma = 1.75$ y $\Delta F_{TH} = 26 - 22 f_{\min}/f_y$ ksi) también cumple:
interior $1\,186.4 \le 1\,620.8$; borde $1\,023.8 \le 1\,593.1\uKg$.}

\paragraph{d) Cortante.} Las losas diseñadas por flexión con el método de franjas de AASHTO 4.6.2.3 se consideran
satisfactorias por corte (AASHTO 5.12.2.1), así que no se diseña refuerzo transversal por corte.

\subsection*{1.12 Acero de distribución y de temperatura}
Refuerzo transversal inferior (distribución), como porcentaje del acero principal:
\begin{align*}
\%\,A_{s,\text{dist}} &= \frac{1750}{\sqrt{S}} = \frac{1750}{\sqrt{6\,096}} = 22.41\% \le 50\%
 \refn{AASHTO 9.7.3.2}\\
A_{s,\text{dist}} &= 22.41\% \times 34.00 = 7.62\uCmm
 \ \Rightarrow\ \boxed{\varnothing\,5/8''\ @\ 25.0\uCm}\quad (8.00\uCmm)
\end{align*}
Se calcula con el mayor $A_s$ colocado (franja de borde) para usar un solo espaciamiento en todo el ancho;
va transversal, en la cara inferior, sobre el acero principal.

Refuerzo superior por temperatura (fórmula de clase, $b = 1000$ mm, $h$ en mm, $f_y$ en MPa):
\begin{align*}
A_{st} &= \frac{0.75\,b\,h}{f_y} = \frac{0.75(1000)(400)}{411.9} = 728.4\ \mathrm{mm^2/m} = 7.28\uCmm\\
 &\Rightarrow\quad \boxed{\varnothing\,1/2''\ @\ 17.5\uCm} \quad(7.37\uCmm)
\end{align*}
Va en la cara superior, en ambas direcciones.

\subsection*{1.13 Gálibo hidráulico}
El fondo de la losa debe quedar por lo menos 2.50 m sobre el nivel de aguas máximas extraordinarias (NAME):
\begin{equation*}
\text{Cota de rasante} \ \ge\ \text{NAME} + 2.50 + t + e_{asf} = \text{NAME} + 2.50 + 0.40 + 0.0762
 = \boxed{\text{NAME} + 2.976\uM}
\end{equation*}
El NAME sale del estudio hidrológico; con él se fija la altura de los estribos.

\subsection*{1.14 Resumen}
\begin{resultado}
\begin{tabularx}{\linewidth}{@{}>{\bfseries}lX@{}}
Vía & Autopista de Segunda Clase; dos puentes gemelos de $W = 15.65$ m, $N_L = 3$\\
Sección & baranda 0.15 + vereda 1.50 + barrera 0.40 + berma 3.00 + calzada 7.20
 + berma 3.00 + barrera 0.40 = 15.65 m\\
Superestructura & losa maciza simplemente apoyada, $L = 6.096$ m, $t = 0.40$ m, $f'_c = 420\uKg$\\
Momentos de diseño & $M_u = 34.363$ (interior), $38.794$ (borde separador) y $14.279$ t$\cdot$m/m (borde vereda)\\
Acero principal inferior & $\varnothing\,1''\ @\ 17.5$ cm (interior y borde separador, 1.518 m)
 \quad $\varnothing\,1''\ @\ 30.0$ cm (borde vereda, 1.636 m)\\
Distribución inferior & $\varnothing\,5/8''\ @\ 25.0$ cm, transversal\\
Temperatura superior & $\varnothing\,1/2''\ @\ 17.5$ cm, en ambas direcciones\\
Rasante & $\ge$ NAME + 2.976 m
\end{tabularx}
\end{resultado}

\subsection*{1.15 Cuadro de verificaciones (demanda / capacidad)}
\begin{center}\small
\begin{tabular}{p{7.4cm}rrlrl}
\toprule
Verificación & Demanda & Capacidad & Unidad & D/C & Estado\\
\midrule
Resistencia I, $M_u \le \phi M_n$ (franja interior) & 34.36 & 37.22 & t$\cdot$m/m & 0.92 & \textcolor{ok}{\textbf{CUMPLE}} \\
Ductilidad $c/d \le 0.42$ (franja interior) & 0.128 & 0.420 &  & 0.30 & \textcolor{ok}{\textbf{CUMPLE}} \\
Acero mínimo $\min(M_{cr};\,1.33M_u) \le \phi M_n$ (franja interior) & 11.79 & 37.22 & t$\cdot$m/m & 0.32 & \textcolor{ok}{\textbf{CUMPLE}} \\
Cuantía mínima $\rho \ge 0.03\,f'_c/f_y$ (franja interior) & 10.65 & 29.14 & cm$^2$/m & 0.37 & \textcolor{ok}{\textbf{CUMPLE}} \\
Servicio I, $f_{ss} \le 0.6 f_y$ (franja interior) & 2\,140.47 & 2\,520.00 & kg/cm$^2$ & 0.85 & \textcolor{ok}{\textbf{CUMPLE}} \\
Fisuración, $s \le s_{\max}$ (franja interior) & 17.50 & 28.21 & cm & 0.62 & \textcolor{ok}{\textbf{CUMPLE}} \\
Fatiga I, $\Delta f \le \Delta F_{TH}$ (franja interior) & 1\,016.91 & 1\,501.68 & kg/cm$^2$ & 0.68 & \textcolor{ok}{\textbf{CUMPLE}} \\
Resistencia I, $M_u \le \phi M_n$ (franja borde separador) & 38.79 & 43.05 & t$\cdot$m/m & 0.90 & \textcolor{ok}{\textbf{CUMPLE}} \\
Ductilidad $c/d \le 0.42$ (franja borde separador) & 0.149 & 0.420 &  & 0.35 & \textcolor{ok}{\textbf{CUMPLE}} \\
Acero mínimo $\min(M_{cr};\,1.33M_u) \le \phi M_n$ (franja borde separador) & 11.79 & 43.05 & t$\cdot$m/m & 0.27 & \textcolor{ok}{\textbf{CUMPLE}} \\
Cuantía mínima $\rho \ge 0.03\,f'_c/f_y$ (franja borde separador) & 10.65 & 34.00 & cm$^2$/m & 0.31 & \textcolor{ok}{\textbf{CUMPLE}} \\
Servicio I, $f_{ss} \le 0.6 f_y$ (franja borde separador) & 2\,115.03 & 2\,520.00 & kg/cm$^2$ & 0.84 & \textcolor{ok}{\textbf{CUMPLE}} \\
Fisuración, $s \le s_{\max}$ (franja borde separador) & 15.00 & 28.66 & cm & 0.52 & \textcolor{ok}{\textbf{CUMPLE}} \\
Fatiga I, $\Delta f \le \Delta F_{TH}$ (franja borde separador) & 877.57 & 1\,476.88 & kg/cm$^2$ & 0.59 & \textcolor{ok}{\textbf{CUMPLE}} \\
Resistencia I, $M_u \le \phi M_n$ (franja borde vereda) & 14.28 & 22.17 & t$\cdot$m/m & 0.64 & \textcolor{ok}{\textbf{CUMPLE}} \\
Ductilidad $c/d \le 0.42$ (franja borde vereda) & 0.075 & 0.420 &  & 0.18 & \textcolor{ok}{\textbf{CUMPLE}} \\
Acero mínimo $\min(M_{cr};\,1.33M_u) \le \phi M_n$ (franja borde vereda) & 11.79 & 22.17 & t$\cdot$m/m & 0.53 & \textcolor{ok}{\textbf{CUMPLE}} \\
Cuantía mínima $\rho \ge 0.03\,f'_c/f_y$ (franja borde vereda) & 10.65 & 17.00 & cm$^2$/m & 0.63 & \textcolor{ok}{\textbf{CUMPLE}} \\
Servicio I, $f_{ss} \le 0.6 f_y$ (franja borde vereda) & 1\,833.64 & 2\,520.00 & kg/cm$^2$ & 0.73 & \textcolor{ok}{\textbf{CUMPLE}} \\
Fisuración, $s \le s_{\max}$ (franja borde vereda) & 30.00 & 34.44 & cm & 0.87 & \textcolor{ok}{\textbf{CUMPLE}} \\
Acero de distribución & 7.62 & 8.00 & cm$^2$/m & 0.95 & \textcolor{ok}{\textbf{CUMPLE}} \\
Acero de temperatura (cara superior) & 7.28 & 7.37 & cm$^2$/m & 0.99 & \textcolor{ok}{\textbf{CUMPLE}} \\
\bottomrule
\end{tabular}
\end{center}

%==========================================================================
\newpage
\section*{PREGUNTA 02 --- Reacción y momento máximo por el Teorema de Barré}

\subsection*{2.1 Datos}
Vehículo de 5 ejes: $7$ --- 3.50 m --- $9$ --- 1.30 m --- $9$ --- 4.50 m --- $9$ --- 1.30 m --- $9$ (t).
Luz $L = 20$ ft $= 6.096\uM$. Largo del vehículo $= 3.50 + 1.30 + 4.50 + 1.30 = 10.60\uM > L$:
\textbf{nunca están todos los ejes sobre el puente}, así que hay que ver qué grupo de ejes cabe en cada posición.

\subsection*{2.2 Reacción máxima}
La reacción $R_A = \sum P_i\,(L - x_i)/L$ varía linealmente mientras no entre ni salga un eje, así que su máximo
ocurre con un eje sobre el apoyo. Se prueba cada eje sobre el apoyo, con el vehículo en ambos sentidos:

\begin{center}\small
\begin{tabular}{cccr}
\toprule
Sentido & Eje sobre el apoyo & Ejes sobre la luz & $R_A$ (t)\\
\midrule
directo & 1 & 1, 2, 3 & 12.746 \\
directo & 2 & 2, 3, 4 & 16.518 \\
directo & 3 & 3, 4, 5 & 11.793 \\
directo & 4 & 4, 5 & 16.081 \\
directo & 5 & 5 & 9.000 \\
inverso & 5 & 3, 4, 5 & 16.518 \\
inverso & 4 & 2, 3, 4 & 11.793 \\
inverso & 3 & 1, 2, 3 & 17.569 \\
inverso & 2 & 1, 2 & 11.981 \\
inverso & 1 & 1 & 7.000 \\
\bottomrule
\end{tabular}
\end{center}

Gobierna el eje 3 sobre el apoyo, con los ejes 1 y 2 sobre la luz:
\begin{equation*}
R_{\max} = 9\cdot\dfrac{6.096-0.00}{6.096} + 9\cdot\dfrac{6.096-1.30}{6.096} + 7\cdot\dfrac{6.096-4.80}{6.096} = \boxed{17.569\uT}
\end{equation*}

\begin{figure}[H]
\centering
\begin{tikzpicture}[x=1.7224cm,y=1cm]
\clip (-1.4,-1.3) rectangle ({6.0960+1.4},2.15);
\vigaSA{6.0960}
\draw[-{Stealth[length=2.2mm]}, very thick, carga] (4.800,1.25) -- (4.800,0.10) node[pos=0, above, font=\scriptsize, text=carga, align=center] {\textbf{1}\\7.00};
\draw[-{Stealth[length=2.2mm]}, very thick, carga] (1.300,1.25) -- (1.300,0.10) node[pos=0, above, font=\scriptsize, text=carga, align=center] {\textbf{2}\\9.00};
\draw[-{Stealth[length=2.2mm]}, very thick, carga] (-0.000,1.25) -- (-0.000,0.10) node[pos=0, above, font=\scriptsize, text=carga, align=center] {\textbf{3}\\9.00};
\draw[-{Stealth[length=2.2mm]}, very thick, gray!55] (-4.500,1.25) -- (-4.500,0.10) node[pos=0, above, font=\scriptsize, text=gray!55, align=center] {\textbf{4}\\9.00};
\draw[-{Stealth[length=2.2mm]}, very thick, gray!55] (-5.800,1.25) -- (-5.800,0.10) node[pos=0, above, font=\scriptsize, text=gray!55, align=center] {\textbf{5}\\9.00};
\draw[<->,>=Stealth] (0,-1.05) -- (6.0960,-1.05) node[midway,fill=white,font=\scriptsize,inner sep=1pt]{$L = 6.096$ m};
\end{tikzpicture}
\caption{Posición de la reacción máxima en el apoyo A (en gris, los ejes fuera de la luz).}
\end{figure}

\subsection*{2.3 Momento máximo (Teorema de Barré)}
Para cada grupo de ejes contiguos que cabe en la luz y cada eje crítico se ubica el centro de la luz en el punto
medio entre el eje y la resultante del grupo ($x = L/2 + e/2$). La posición es válida solo si en ella están
sobre la luz exactamente los ejes del grupo:

\begin{center}\small
\begin{tabular}{ccccccc}
\toprule
Grupo & Eje crítico & $R$ (t) & $e$ (m) & $x$ crítico (m) & ¿Posición válida? & $M$ (t$\cdot$m)\\
\midrule
1--2 & 1 & 16.00 & -1.969 & 2.064 & \textcolor{ok}{sí} & 11.177 \\
1--2 & 2 & 16.00 & 1.531 & 3.814 & \textcolor{mal}{no} & -- \\
1--2--3 & 1 & 25.00 & -2.988 & 1.554 & \textcolor{mal}{no} & -- \\
1--2--3 & 2 & 25.00 & 0.512 & 3.304 & \textcolor{mal}{no} & -- \\
1--2--3 & 3 & 25.00 & 1.812 & 3.954 & \textcolor{mal}{no} & -- \\
2--3 & 2 & 18.00 & -0.650 & 2.723 & \textcolor{ok}{sí} & 21.894 \\
2--3 & 3 & 18.00 & 0.650 & 3.373 & \textcolor{ok}{sí} & 21.894 \\
2--3--4 & 2 & 27.00 & -2.367 & 1.865 & \textcolor{mal}{no} & -- \\
2--3--4 & 3 & 27.00 & -1.067 & 2.515 & \textcolor{mal}{no} & -- \\
2--3--4 & 4 & 27.00 & 3.433 & 4.765 & \textcolor{mal}{no} & -- \\
3--4 & 3 & 18.00 & -2.250 & 1.923 & \textcolor{mal}{no} & -- \\
3--4 & 4 & 18.00 & 2.250 & 4.173 & \textcolor{mal}{no} & -- \\
3--4--5 & 3 & 27.00 & -3.433 & 1.331 & \textcolor{mal}{no} & -- \\
3--4--5 & 4 & 27.00 & 1.067 & 3.581 & \textcolor{mal}{no} & -- \\
3--4--5 & 5 & 27.00 & 2.367 & 4.231 & \textcolor{mal}{no} & -- \\
4--5 & 4 & 18.00 & -0.650 & 2.723 & \textcolor{ok}{sí} & 21.894 \\
4--5 & 5 & 18.00 & 0.650 & 3.373 & \textcolor{ok}{sí} & 21.894 \\
\bottomrule
\end{tabular}
\end{center}

Los grupos de tres ejes no son válidos: al ubicarlos según Barré, un eje del grupo queda fuera o entra otro eje.
Las posiciones válidas son los ejes dobles 2--3 y 4--5 (9 t + 9 t a 1.30 m):
\begin{equation*}
M_{\max} = \frac{2P}{L}\left(\frac{L}{2} - \frac{s}{4}\right)^{2}
 = \frac{2(9)}{6.096}\left(\frac{6.096}{2} - \frac{1.30}{4}\right)^{2} = \boxed{21.894\uTm}
\end{equation*}
bajo el eje 2, a $x = 2.723\uM$ del apoyo (0.325 m del centro de la luz).

\textbf{Reacciones en la posición de Barré.} La resultante de los ejes 2 y 3 ($R = 18.00\uT$) queda a
$\bar{x} = 3.373\uM$ del apoyo A:
\begin{align*}
R_A &= \frac{R\,(L - \bar{x})}{L} = \frac{18.00\,(6.096 - 3.373)}{6.096} = 8.040\uT \qquad
R_B = R - R_A = 18.00 - 8.040 = 9.960\uT\\
M_{\max} &= R_A \cdot x = 8.040 \times 2.723 = 21.894\uTm \ok
\end{align*}
Comparación: un eje de 9 t solo al centro da $PL/4 = 13.716\uTm$, y con el eje 2 al centro (el eje 3
también sobre la luz) se obtiene $21.582\uTm$; ambos son menores. Un barrido numérico del vehículo
cada 1 mm confirma $M_{\max} = 21.894\uTm$.

\begin{figure}[H]
\centering
\begin{tikzpicture}[x=1.7224cm,y=1cm]
\clip (-1.4,-1.65) rectangle ({6.0960+1.4},2.15);
\vigaSA{6.0960}
\draw[-{Stealth[length=2.2mm]}, very thick, gray!55] (-0.777,1.25) -- (-0.777,0.10) node[pos=0, above, font=\scriptsize, text=gray!55, align=center] {\textbf{1}\\7.00};
\draw[-{Stealth[length=2.2mm]}, very thick, carga] (2.723,1.25) -- (2.723,0.10) node[pos=0, above, font=\scriptsize, text=carga, align=center] {\textbf{2}\\9.00};
\draw[-{Stealth[length=2.2mm]}, very thick, carga] (4.023,1.25) -- (4.023,0.10) node[pos=0, above, font=\scriptsize, text=carga, align=center] {\textbf{3}\\9.00};
\draw[-{Stealth[length=2.2mm]}, very thick, gray!55] (8.523,1.25) -- (8.523,0.10) node[pos=0, above, font=\scriptsize, text=gray!55, align=center] {\textbf{4}\\9.00};
\draw[-{Stealth[length=2.2mm]}, very thick, gray!55] (9.823,1.25) -- (9.823,0.10) node[pos=0, above, font=\scriptsize, text=gray!55, align=center] {\textbf{5}\\9.00};
\draw[dashed,gray] ({6.0960/2},-0.25) -- ({6.0960/2},1.95) node[above,font=\scriptsize]{$L/2$};
\draw[<->,>=Stealth] (0,-1.0) -- (2.7230,-1.0) node[midway,below,font=\scriptsize]{$x = 2.723$ m};
\draw[<->,>=Stealth] (0,-1.45) -- (6.0960,-1.45) node[midway,fill=white,font=\scriptsize,inner sep=1pt]{$L = 6.096$ m};
\end{tikzpicture}
\caption{Posición de Barré para el momento máximo (eje doble 2--3).}
\end{figure}

\begin{resultado}
\textbf{Respuestas, Pregunta 02:} \quad $R_{\max} = 17.569\uT$ \qquad $M_{\max} = 21.894\uTm$\\[2pt]
{\small En la posición de Barré: $R_A = 8.040\uT$ y $R_B = 9.960\uT$.}
\end{resultado}

\vspace{4pt}
{\small\textit{Cálculos generados con \texttt{puente\_losa.py} (mismos datos que este documento). Supuestos por
confirmar: espesor de vereda de 0.20 m (1.2), factor de exposición $\gamma_e$ (1.11-b) y edición
de AASHTO para Fatiga (1.11-c).}}

\end{document}
